\title{Large Language Models Can Automatically Engineer Features for Few-Shot Tabular Learning}

\begin{abstract}
Large Language Models (LLMs) have demonstrated exceptional capability in addressing complex and novel reasoning tasks, positioning them as a promising resource for tabular learning, which is central to numerous practical applications. In this work, we introduce a new in-context learning approach called \textsf{FeatLLM}, wherein LLMs are leveraged to design features that optimize the data representation for tabular prediction problems. The features crafted by the LLM serve as the basis for predicting class probabilities through a straightforward downstream machine learning method, such as linear regression, enabling robust performance in few-shot learning scenarios. Importantly, at prediction time, the \textsf{FeatLLM} approach employs only the derived features together with a basic predictive model, dispensing with recurrent LLM queries for each test instance, as required by prior LLM-centric methods. Additionally, \textsf{FeatLLM} only necessitates API-level interaction with LLMs, thus addressing constraints imposed by prompt length. Through extensive evaluation on diverse tabular datasets spanning multiple fields, \textsf{FeatLLM} is shown to generate high-quality decision rules, achieving substantially better results—surpassing competitors like TabLLM and STUNT by 10\% on average.
\end{abstract}

\section{Introduction}
LLMs have recently achieved remarkable success in generating appropriate outputs for tasks unencountered during training, all without the need for dedicated fine-tuning~\cite{creswell2022selection,imani2023mathprompter,taylor2022galactica}. Their training on vast textual datasets imparts broad prior knowledge, enabling LLMs to replace domain experts in various scenarios~\cite{Genomes2021,singhal2023large,wilcox2003role} and aiding the automation of tasks such as conversational analysis~\cite{finch2023leveraging} and automated program repair~\cite{fan2023automated}. By including task descriptions and a handful of relevant examples in prompt construction, LLMs can infer the context, prioritize crucial features, and reference illustrative instances~\cite{brown2020language,zhao2023survey}. This attests to LLMs’ proficiency in data interpretation and underscores their potential to function as skilled feature engineers.

Extending the application of LLMs' prior knowledge and inferential abilities, tabular learning has begun to benefit from these advances. As an example, information regarding the increased susceptibility of certain ailments among elderly individuals assists disease prediction tasks. Recent studies have operationalized this by formulating tasks and feature explanations in natural language, transforming the data into text, and submitting it to the LLM~\cite{dinh2022lift,wang2023anypredict}. This approach is frequently combined with techniques such as in-context learning~\cite{nam2023semi} or methods for efficient fine-tuning~\cite{dinh2022lift,hegselmann2023tabllm}. The implicit knowledge accessible via LLMs has been found to be highly effective, providing superior performance compared to standard tabular learning baselines in scenarios with limited labeled examples~\cite{hegselmann2023tabllm}.

Existing LLM-based tabular learning techniques encounter several significant hurdles. To generate end-to-end predictions, these methods necessitate conducting at least one LLM inference for every individual data sample, leading to substantial computational demands. Achieving satisfactory predictive performance typically depends on fine-tuning the LLM, but this is impractical when LLMs lack public training codebases or accessible fine-tuning APIs—a concern heightened by the fact that many state-of-the-art LLMs restrict user interaction to inference APIs only~\cite{anil2023palm,openai2023gpt4}. Another complication arises when tabular datasets have numerous features; as the prompt text length increases, the applicability and effectiveness of existing approaches diminish or become unmanageable~\cite{liu2023lost}. Collectively, these issues hinder the practical implementation of LLM-driven tabular learning in real-world settings.

In response, we advocate for an alternative paradigm that repositions LLMs as tools for feature engineering, moving away from direct prediction tasks. We present a novel in-context learning methodology, \textsf{FeatLLM}, which is designed to discern and articulate the “criteria” guiding LLM-generated predictions. For example, in disease classification scenarios, the LLM can enumerate the specific rules that associate feature patterns with disease identification. These derived rules are subsequently transformed into new features, supplanting original inputs for downstream modeling tasks. This shift delegates feature creation to the LLM, enabling the use of simpler models thereafter—thus reducing inference time, since complex models are no longer necessary. Crucially, this framework leverages the LLM's in-context learning, allowing feature extraction without training the LLM: informative examples can simply be incorporated into the prompts. Additionally, techniques such as feature bagging from ensemble learning literature~\cite{breiman2001random} offer a strategy to control the prompt length, helping mitigate the challenge posed by large numbers of features.

\textsf{FeatLLM} organizes the process of feature creation within the prompt into two distinct stages: first, comprehending the underlying problem and identifying how features are connected to target variables; second, leveraging insights gained along with few-shot samples to formulate clear and discriminative rules for class separation. This methodological reasoning encourages LLMs to exclude irrelevant features during rule generation. Subsequently, the extracted rules are implemented directly on the dataset (as program code), converting each sample into a binary indicator corresponding to rule compliance. A simple predictive model is then trained on these binary features to estimate class probabilities. To enhance robustness, the method employs an ensemble technique by repeatedly producing feature sets through feature bagging. Adjusting both the number of features and sample size during bagging effectively addresses the issue of prompt size limitations. \textsf{FeatLLM} operates without any training requirements and incurs minimal inference expense, making LLM deployment practical across a wide range of applications.

Our experiments with \textsf{FeatLLM} span 13 diverse tabular datasets under limited-sample conditions, demonstrating its consistently strong and reliable performance. We illustrate that LLMs effectively harness prior knowledge as well as information contained within the data, yielding superior feature extraction. The framework achieves better results than state-of-the-art few-shot learning baselines in numerous scenarios. Code for our approach is available via the anonymized GitHub repository at~\texttt{https://github.com/Sungwon-Han/FeatLLM}.

\section{Related Work}

\subsection{Few-Shot Learning with Tabular Data} Recent developments in few-shot learning have facilitated the acquisition of robust, transferable representations from datasets with limited labeled examples~\cite{chen2018closer}. Although pioneering works predominantly addressed image domains, few-shot techniques have expanded to encompass diverse modalities such as text, audio, and, increasingly, tabular datasets~\cite{majumder2022few,nam2023stunt,schick2021exploiting}. When models are trained on scant labeled data, there is a heightened risk of inadvertently capturing misleading correlations~\cite{bartlett2021deep}. As a result, contemporary research has incorporated supplementary information sources or integrated domain-specific priors to enforce beneficial inductive biases during learning. Strategies include leveraging large volumes of unlabeled data via data augmentation in semi-supervised settings~\cite{ucar2021subtab,yoon2020vime}, and pursuing unsupervised pre-training for improved model initialization that is agnostic to the specific task~\cite{somepalli2021saint}. Common unsupervised objectives feature mechanisms such as data masking or corruption, followed by reconstructive learning~\cite{arik2021tabnet,majmundar2022met}, or leveraging augmentation for contrastive learning~\cite{bahri2022scarf,chen2020simple}. Nevertheless, the complexity and heterogeneity of tabular data hampers the development of universal augmentation strategies, limiting their efficacy for few-shot scenarios~\cite{nam2023stunt}. 

To address this, recent research has advocated for training on a wide array of tabular datasets not restricted to the task-specific domain, followed by transfer learning to adapt to the target task~\cite{zhang2023generative,levin2022transfer,zhu2023xtab}. For example, TransTab~\cite{wang2022transtab} utilizes transformer architectures that discern semantic relationships between columns using column names and table contents. The accumulated knowledge across numerous tables facilitates both zero-shot and few-shot inference for novel tasks. Conversely, TabPFN~\cite{hollmann2022tabpfn} generates synthetic datasets from diverse distribution types for pre-training Bayesian neural networks. Inference after pre-training involves providing both the target task’s training and test examples as input without additional optimization. This multi-source prior enables these models to outperform conventional tree-based methods and demonstrates notable advantages in few-shot contexts.

\subsection{Language-Interfaced Tabular Learning} Integrating large language models (LLMs)~\cite{anil2023palm,openai2023gpt4} offers another promising avenue for injecting prior knowledge into tabular learning. By virtue of their extensive training on diverse textual data, LLMs excel in generalizing to novel tasks, establishing their relevance across multiple domains~\cite{brown2020language}. Recent efforts have leveraged the inherent knowledge in LLMs for tabular data problems, commonly transforming tabular datasets into serialized text formats and embedding them within LLM prompt inputs~\cite{dinh2022lift,hegselmann2023tabllm,wang2023anypredict}. Prompt construction typically integrates tabular values, feature metadata, and task descriptions, enabling the model to reason about task-feature relationships for downstream inference. Approaches in this area span specialized parameter-efficient fine-tuning methods such as LoRA~\cite{hu2021lora} and IA3~\cite{liu2022few}, as well as prompt-based in-context learning, wherein a handful of exemplary cases are provided within the input—without any additional model training~\cite{dinh2022lift,hegselmann2023tabllm,nam2023semi,slack2023tablet}.

While existing approaches have demonstrated notable success in advancing few-shot tabular learning, their reliance on routing every inference through the LLM presents practical difficulties, particularly for industrial settings. In contrast to the traditional black-box framework that predicts each data point using an end-to-end LLM pipeline, this work shifts focus to having the LLM directly deduce the class-specific conditions or rules. Through \textsf{FeatLLM}, these derived rules are converted into programmatic code, yielding supplementary features that enable prediction independent of LLM inference, while simultaneously drawing on in-context learning to distill rules from existing knowledge and the information present in few-shot examples.

\section{Methods}
\textsf{FeatLLM} is structured to extract "rules"—the defining criteria for each class—from limited input samples, as illustrated in Figure~\ref{fig:main_model}. The rules generated by the LLM are then interpreted to construct new binary features for every sample, with each feature representing whether the corresponding rule holds for that sample. These binary indicators are subsequently incorporated in a dot product calculation with non-negative, trainable weights, which yields an estimate for the class likelihood. The training process learns both the relative importance of each feature and their influence on predicting the target class (see Section~\ref{Sec:training}). To enhance robustness, this procedure is performed repeatedly with bagging, followed by ensemble aggregation. The following sections detail the problem setup and the specifics of each step.

\vspace*{-0.12in}
\paragraph{Problem formulation.} Consider a tabular dataset comprising $N$ labeled instances, $\mathcal{D} = \{(\mathbf{x}^i, \mathbf{y}^i)\}_{i=1}^{N}$. Each data sample $\mathbf{x}^i$ represents a $d$-dimensional vector corresponding to $d$ distinct input features. These features are annotated with descriptive natural language names, collected as $F = \{f_j\}_{j=1}^{d}$, and their explanations are documented separately. Within a classification context, the label $\mathbf{y}^i$ assigns each instance to one of several predefined categories. For $k$-shot learning scenarios, a subset of $k$ ($k < N$) randomly chosen labeled examples is used to train the model.

\subsection{Prompt Design for Extracting Rules}
\label{Sec:prompt_design}
To facilitate rule extraction from LLMs using informed reasoning procedures, we orient the problem-solving pathway to reflect the strategy an expert human would adopt when handling tabular data tasks. The constructed prompt consists of three primary sections, as illustrated in Fig.~\ref{fig:prompt_for_rule} and detailed in Appendix~\ref{sec:example_rule}:

\vspace*{-0.12in}
\paragraph{Basic information description.} The initial component supplies crucial context necessary for addressing the task (refer to orange highlighted text in Fig.~\ref{fig:prompt_for_rule}). This includes a description of the specific task—paired with label information—feature summaries, and example instances. The task is articulated as a direct question, such as ``Does this patient's myocardial infarction complications data indicate chronic heart failure? Yes or no?". Additional examples can be found in Appendix~\ref{sec:task_desc}. Feature descriptions specify the data type for each feature and may optionally provide a definition. Categorical features can be accompanied by illustrative category examples. For demonstration purposes, a limited number of training cases are converted to textual form, together with their true labels. Serialization adheres to conventions established by prior work~\cite{dinh2022lift,hegselmann2023tabllm}:

\begin{align}
&\text{Serialize}(\mathbf{x}^i,\mathbf{y}^i, F) = \nonumber \\ 
&``\ f_1\ \text{is}\ \mathbf{x}^i_1.\ ... \ f_d\  \text{is}\  \mathbf{x}^i_d.\ \ \text{Answer:}\  \mathbf{y}^i\ ”
\end{align}

\vspace*{-0.12in}
\paragraph{Reasoning instruction.} Our approach focuses on improving the LLM's reasoning capability by incorporating explicit reasoning instructions (highlighted in blue in Fig.~\ref{fig:prompt_for_rule}). The reasoning instructions start with an introductory statement, echoing the chain-of-thought strategy~\cite{wei2022chain}. Next, we organize the model’s inference into two distinct phases. Initially, the LLM is prompted to hypothesize causal connections or behavioral tendencies linking the features to the task description, utilizing its own general understanding or common sense without relying on provided examples. This ensures that the LLM leverages its prior knowledge effectively. Subsequently, in the second phase, the LLM integrates insights from the example demonstrations along with outcomes from the first step to generate class-specific rules. This staged approach helps steer the model away from picking up on coincidental associations with irrelevant features and instead emphasizes those attributes that are most pertinent. To balance the number of rules generated—thus avoiding both sparse rule sets (underfitting) and overly voluminous prompts—we fix the number of extracted rules per class at 10\footnote{Further discussion regarding rule count is provided in Section~\ref{sec:analysis}.}. Moreover, while generating rules, the LLM references feature type information from the descriptive portion, ensuring rule compatibility and preventing mismatches.

\vspace*{-0.12in}
\paragraph{Response instruction.} To streamline the parsing and application of inferred rules, we specify detailed guidance for the LLM on response formatting (as illustrated in yellow in Fig.~\ref{fig:prompt_for_rule}). This prompt includes explicit directives on how to organize responses for each answer class (response format), as well as how each rule should be represented (rule format). By following these guidelines, the LLM can appropriately format rules according to feature type. If provided without such instructions, the LLM may naturally combine multiple rules by utilizing logical connectors like AND or OR. An example of this output can be found in Appendix~\ref{sec:outcome-rule}.

\subsection{Inferring Class Likelihood via Rules}
\label{Sec:training}

\paragraph{Rule parsing for feature extraction.} The previously constructed rules are employed to generate novel binary features. Each feature corresponds to a specific class, marking whether the sample fulfills the criteria defined by the rules of that class (yielding either ‘0’ or ‘1’). These binary indicators provide information relevant to determining a sample's class membership. Since these rules originate from natural language outputs produced by an LLM, an automated parsing approach must be formulated to enable feature generation. Although we enforce response formatting constraints during rule creation for smoother parsing, LLM outputs can still present unpredictable syntactic variations~\cite{kovavc2023large}. For example, an LLM might rephrase operators, such as replacing ``$>$” and ``$<$” with ``greater than" or ``less than,” introducing parsing complexity.

To mitigate issues caused by parsing noisy textual content, rather than designing elaborate code to handle such cases, we exploit the LLM's capabilities. Owing to its proficiency in capturing semantic meanings in diverse syntactic forms and its aptitude for generating straightforward code (as a result of training on code-rich datasets), we harness these strengths. Our prompt formulation incorporates the target function’s name, clear descriptions of the inputs and outputs, as well as the extracted rules, before passing this prompt to the LLM. Illustrations of these prompts and their outcomes are depicted in Figures~\ref{fig:prompt_for_parsing} and ~\ref{sec:example_parsing}-\ref{sec:example_parse_outcome} in Appendix. The LLM-generated Python code is subsequently run through the \texttt{exec()} function, facilitating the conversion of data.

\vspace*{-0.12in}
\paragraph{Estimating class probabilities.} With access to binary features derived from class-specific rules, an intuitive approach to gauge a sample’s probability of belonging to a class involves tallying the number of rules satisfied for each class (computing the sum across relevant binary features). Nevertheless, some rules and their corresponding features may be more informative than others; thus, their relative importance must be determined using training data. We set out to quantify this by fitting a standard machine learning (ML) model on each class’s set of binary features. A linear model with zero bias is a simple candidate. Representing the binary feature vector for sample $\mathbf{x}^i$ and class $k$ as $\mathbf{z}_k^i \in \{0, 1\}^R$, where $R$ is the rule count per class and $c$ denotes the number of classes, the probability distribution $\mathbf{p}^i$ for the classes can be calculated as:

\begin{align}
\text{logit}_k^i &= \max(\mathbf{w}_k, 0) \cdot \mathbf{z}_k^i, \nonumber \\
\mathbf{p}^i &= \text{Softmax}({[\text{logit}_{1}^i, ..., \text{logit}_{c}^i]}). \label{eq:infer_prob}
\end{align}

In this context, $\mathbf{w}_k$ denotes the learnable parameters within the linear model, serving as indicators for the significance of each corresponding feature. By constraining these weights to reside within a positive subspace, we guarantee that features produced by the LLM do not adversely affect the predicted probability for any class during training. Subsequently, the logit outputs are translated into class probabilities via the softmax transformation. We leverage a few-shot set of training instances to optimize the cross-entropy objective and infer $\mathbf{w}_k$.

\vspace*{-0.12in}
\paragraph{Ensembling with bagging.} The process is iterated multiple times, thereby yielding several distinct inference models whose predictions are aggregated through ensembling. For $T$ runs producing trained weight vectors $\{\mathbf{w}[t]\}_{t=1}^T$, the final output is derived by averaging per-class probabilities $\mathbf{p}^i$, as computed using each $\mathbf{w}[t]$ according to Eq.~\ref{eq:infer_prob}.

To facilitate diversity among the generated rules for ensembling, LLM inference is conducted with a sufficiently elevated temperature. Additionally, we permute the sequence of few-shot examples, since it has been shown that their order can influence LLM responses~\cite{lu2022fantastically}\footnote{A comprehensive assessment of rule diversity is available in Appendix~\ref{sec:diversity}}. Bagging further enhances the variety and reliability of predictions by selecting different subsets of features or data samples for each model instantiation, particularly valuable when dealing with numerous features or sizable training data. The ensemble method confers two principal benefits. Firstly, when some LLM-generated rules prove to be erroneous, accurate rules from alternate trials can offset these mistakes, thereby increasing the framework's robustness; this aligns with the concept of LLM self-consistency~\cite{wang2022self}, as frequently recurring rules across runs are more likely valid. Secondly, the ensembling strategy mitigates constraints posed by the prompt length limit of LLMs. If tabular datasets exceed the permissible prompt size, the bagging approach ensures manageable input sizes while preserving reliable performance. Although any individual rule set might overlook certain feature relationships, an ensemble of diverse weak models compiled from various iterations diminishes the effect of such omissions in single trials.

\section{Experiments}
We assess \textsf{FeatLLM} across a range of tabular datasets and carry out several analyses to explore the mechanisms underlying our framework. Further experimental details, such as tests with alternative LLMs and in-depth qualitative studies, are provided in the Appendix due to space limitations.

\subsection{Performance Evaluation}
\paragraph{Datasets.}
For our study, we use 13 distinct datasets intended for binary and multi-class classification problems: (1) Adult~\cite{asuncion2007uci}, which involves forecasting if an individual's annual income surpasses \$50,000; (2) Bank~\cite{moro2014data}, targeting the prediction of customer subscription to term deposits; (3) Blood~\cite{yeh2009knowledge}, designed to predict the likelihood of donor return; (4) Car~\cite{kadra2021well}, focusing on car quality assessment; (5) Communities~\cite{misc_communities_and_crime_183}, used to estimate the crime rate in a certain area; (6) Credit-g~\cite{kadra2021well}, where the task is to assess whether someone represents a good or bad credit risk; (7) Diabetes\footnote{\texttt{kaggle.com/datasets/uciml/pima-indians-diabetes-database}}, for identifying diabetes in an individual; (8) Heart\footnote{\texttt{kaggle.com/datasets/fedesoriano/heart-failure-prediction}}, predicting coronary artery disease occurrence; and (9) Myocardial~\cite{misc_myocardial_infarction_complications_579}, used to detect chronic heart failure cases.

Additionally, we incorporate recently released and synthetic datasets that have received minimal coverage and were not included during the LLMs' pre-training phase: (10) Cultivars, which predicts whether a soybean cultivar produces a high grain yield\footnote{\texttt{archive.ics.uci.edu/dataset/913}}; (11) NHANES, which classifies a person's record into age groups\footnote{\texttt{archive.ics.uci.edu/dataset/887}}; (12) Sequence-type, a dataset assigning a sequence to one of four categories: arithmetic, geometric, Fibonacci, or Collatz; (13) Solution-mix, which determines if a mixture’s concentration (from four constituent solutions) exceeds 0.5. The first two datasets were added to the UCI Machine Learning Repository post-September 2023, thereby ensuring their exclusion from the training data of the LLM API (gpt-3.5-turbo-0613) used in our experiments. The last two are artificial datasets, further preventing any possibility of memorization by the LLM API during pretraining, thus allowing genuine evaluation.

A summary of the datasets' scale and structural complexity is given in Table~\ref{Tab:basic-desc}.
Each dataset includes explicit naming and description for all features. Notably, datasets such as `Communities' and `Myocardial' contain over 100 attributes, creating substantial challenges regarding the restricted context windows of LLMs.

\vspace*{-0.12in}
\paragraph{Baselines.}
We benchmark our model against nine different baseline methods. The first group of baselines employs traditional approaches for tabular data, including (1) Logistic regression (LogReg), (2) XGBoost~\cite{chen2016xgboost}, and (3) RandomForest~\cite{ho1995random}. Two additional baselines operate under the assumption of having access to large quantities of unlabeled data: (4) SCARF~\cite{bahri2022scarf} and (5) STUNT~\cite{nam2023stunt}; however, collecting such unlabeled datasets can be difficult in practical scenarios. We also include (6) TabPFN~\cite{hollmann2022tabpfn}, which trains using synthetically generated datasets of diverse distributions. The remaining baselines are based on LLM methodologies, consisting of (7) In-context learning (In-context)~\cite{wei2022emergent}, (8) TABLET~\cite{slack2023tablet}, and (9) TabLLM~\cite{hegselmann2023tabllm}. In-context learning approaches append few-shot training samples directly within prompts, eschewing any fine-tuning. TABLET supplements prompts with additional contextual hints, drawn from an external classifier using rule sets and prototypes, to boost inference quality. TabLLM leverages parameter-efficient tuning (like IA3~\cite{liu2022few}) to adapt LLMs using limited training examples.

\vspace*{-0.12in}
\paragraph{Implementation details.}
\textsf{FeatLLM} adopts GPT-3.5 as the primary LLM, although its architecture is flexible and not tied to any specific LLM (refer to Appendix~\ref{sec:backbones} for comparative backbone performance). The LLM's inference temperature is configured at 0.5, while the top-p parameter maintains the API's default value of 1. For feature extraction, 20 ensembles and 10 rules are established. Further analysis regarding hyper-parameter selection and influence can be found in Figure~\ref{fig:hyperparameter2} and Appendix~\ref{sec:hyper-parameter}. The linear model is optimized using Adam with a learning rate of 0.01, trained across 200 epochs. To determine the optimal epoch, $k$-fold cross-validation is utilized. Baseline comparisons employ GPT-3.5 for in-context learning; for TabLLM, T0~\cite{sanh2022multitask} is implemented to mirror the original experimental setup. It is worth noting TabLLM confines LLM usage to models with open checkpoints, thereby excluding GPT-3.5 from consideration. Parameter tuning for conventional machine learning techniques like LogReg, XGBoost, and RandomForest is accomplished via Grid-search coupled with $k$-fold cross-validation. Here, $k$ takes the value 2 or 4, ensuring the training subset contains examples from every class. Appendix~\ref{sec:baselines} contains further implementation specifics.

\vspace*{-0.12in}
\paragraph{Results.}
Tables~\ref{tab:main1} and \ref{tab:main2} display comparative performance results across multiple datasets. Each experimental setup was executed three times, varying the random partitioning of training and testing sets; for each configuration, we report the mean and standard deviation of the resulting AUC scores. Across benchmarks, \textsf{FeatLLM} frequently emerges either as the leading model or as the runner-up when measured against other baselines. Remarkably, our method, with only 4 shots, matches the performance of traditional tabular baselines which utilize 32 shots, as illustrated in Fig.~\ref{fig:compares}. While STUNT and TabPFN yielded remarkable gains for some datasets, they did not deliver consistent, substantial improvements relative to standard machine learning approaches. Further, LLM-based approaches, including in-context learning, TABLET, and TabLLM, were incapable of scaling to tabular data with high feature dimensionality. In contrast, our architecture—leveraging both bagging and ensemble methodologies—successfully maintained strong results even in high-feature settings. The modular ensemble and bagging strategies in our framework could theoretically be transferred to other LLM-powered techniques; however, such application would double per-sample inference costs, which renders this option computationally prohibitive.

\subsection{Analysis \& Discussion}
\label{sec:analysis}
\paragraph{Ablation study.}
To assess the role of individual components within our approach, we conduct a series of ablation experiments targeting: (1) \textsf{-Tuning}, where the weight adjustment phase in the linear model is omitted and logits are calculated through direct summation of binary feature values; (2) \textsf{-Ensemble}, where ensemble aggregation is excluded; (3) \textsf{-Description}, where the feature description is removed; and (4) \textsf{-Reasoning}, where the Step-1 reasoning instruction for rule generation is bypassed.

Table~\ref{tab:ablation} presents the average AUC variation caused by each ablation across datasets. Eliminating any individual component leads to a reduction in model performance. Interestingly, the absence of weight tuning exerts a larger negative impact as the shot count increases, implying that with more data, finely tuned rule weighting is key to improved outcomes. Conversely, both the feature description and the reasoning instruction play a pivotal role in low-shot scenarios, highlighting the importance of efficiently leveraging LLMs' prior knowledge for better performance. Finally, applying the ensemble method as proposed robustly boosts predictive accuracy in every evaluated configuration.

\vspace*{-0.12in}
\paragraph{Training \& inference time comparison.} We evaluate various methods based on their training and inference durations, using the Adult dataset and a training set of 16 samples. All models are run on a single A100 GPU by default, except TabLLM, which utilizes four A100 GPUs to enable model parallelism. Table~\ref{tab:cost} displays the runtime results. \textsf{FeatLLM} demonstrates efficient inference performance, closely matching that of traditional techniques, whereas LLM-based alternatives—including In-context learning, TABLET, and TabLLM—incur significantly greater inference times. As for training, \textsf{FeatLLM} achieves similar speeds to other LLM-based approaches; it operates with just 30 API calls, a manageable expense that supports parallel execution.

\vspace*{-0.12in}
\paragraph{Quality of features generated by \textsf{FeatLLM}.}
To assess the utility of rule-based features produced by \textsf{FeatLLM} in practical applications, we carry out an experiment that calculates the average AUROC between these features and the target labels. We then determine the proportion of features per dataset with AUROC values exceeding 0.5. The findings are reported in Table~\ref{tab:quality}. These results indicate that most of the constructed rules are aligned with the target variable and contribute meaningfully to the prediction task (see ``Before tuning" column). Furthermore, when tuning the linear model, features that prove uninformative—characterized by learned weights below a designated threshold—are systematically excluded (see ``After tuning" column). Here, the threshold for feature relevance is established at 0.1, guided by the distribution of learned weights.

\vspace*{-0.12in}
\paragraph{Impact of incorporating spurious correlations.} To evaluate how well different techniques mitigate the influence of irrelevant spurious correlations under few-shot conditions, we performed a targeted analysis. In our experiment, we augmented the Heart dataset by appending randomly selected columns from the Adult dataset that are not pertinent to the Heart prediction task. We assessed the models’ performance as the number of these extraneous columns increased incrementally. To ensure an equitable comparison, all models were given 8 labeled samples (shots), and the setup did not include any unlabeled data, excluding approaches such as SCARF and STUNT from the analysis.

Figure~\ref{fig:spurious} presents the observed performance trajectories affected by spurious correlations. Among the evaluated methods, \textsf{FeatLLM} exhibited the minimal decline in accuracy when exposed to unrelated features. This resilience is likely attributable to the design of our framework, which explicitly emphasizes the importance of the connection between the prediction task and input features using prior knowledge during rule extraction; as a result, rule formation based on irrelevant columns is curtailed, supporting robustness. Empirically, rules grounded in noisy columns comprised merely 1-2\% of all rules. Nonetheless, it is also apparent that \textsf{FeatLLM} may inadvertently capture spurious associations or misinterpret causal linkages when columns derived from semantically related datasets are appended arbitrarily (refer to Appendix~\ref{sec:spurious-appendix} for details).

\vspace*{-0.12in}
\paragraph{Hyper-parameter analysis}
\textsf{FeatLLM} incorporates two critical hyper-parameters: the size of the ensemble and the quantity of rules extracted per class during LLM queries. To evaluate how these settings influence model efficacy, we performed a series of experiments. Our findings indicate that increasing the ensemble size initially yields notable gains in performance, but these improvements plateau around an ensemble size of 20 (refer to Fig.~\ref{fig:appendix-hyperparameter1} in Appendix). Regarding the number of rules, we observed no significant statistical variation in outcomes; however, selecting 10 rules presents an optimal balance, minimizing prompt length while maintaining strong performance (see Fig.~\ref{fig:hyperparameter2}). We attribute this result to the fact that increasing the rule count expands input dimensionality, which amplifies the information content, but also introduces a risk of overfitting, especially in the context of limited samples.

\section{Conclusion}
This work introduces \textsf{FeatLLM}, a novel method that pushes forward the boundaries in LLM-driven few-shot tabular learning. Unlike conventional approaches where LLMs provide sample-specific predictions, \textsf{FeatLLM} leverages LLMs as feature engineers to generate prediction criteria. By adopting a rule generation strategy and an ensemble framework based on bagging, \textsf{FeatLLM} substantially reduces inference costs, requires only API-level integration rather than extensive training, and is not restricted by data feature dimensionality. These innovations facilitate robust predictive performance and establish \textsf{FeatLLM} as a compelling and efficient solution for real-world tabular data challenges.

\textsf{FeatLLM} has been purposely architected for the low-shot learning setting. Future work will involve extending its functionalities to datasets with larger sample sizes, investigating alternative feature generation strategies beyond rule-based methods, and enhancing interpretability to broaden its applicability across different domains.

\section{Broader Impact} The framework, \textsf{FeatLLM}, seeks to capitalize on the prior knowledge encoded within LLMs for tabular learning, aiming to surpass conventional supervised approaches. By improving data efficiency, our approach effectively addresses the challenge of overfitting in environments where training samples are scarce, supplementing limited labeled data with LLM-informed insights. This has substantial practical relevance, especially for domains such as finance and healthcare where obtaining labeled data is costly or difficult, and pretrained LLMs often contain useful domain-specific information. Moving forward, an essential area for investigation involves assessing how societal biases and misinformation inherent in the prior knowledge of LLMs may influence outcomes. Since the framework directly integrates these priors into its predictions, such biases can potentially supersede the signal from observed or curated labeled data, underscoring the importance of careful analysis for responsible adoption.

\end{document}